\documentclass[a4paper,11pt]{article}
\usepackage[utf8]{inputenc}
\usepackage[T1]{fontenc}
\usepackage[polish]{babel}
\usepackage{amsmath, amssymb, amsthm}
\usepackage{algorithm}
\usepackage{algorithmic}
\usepackage{geometry}
\usepackage{tikz}
\usepackage{array}
\usepackage{xcolor}
\geometry{margin=2.5cm}

\theoremstyle{definition}
\newtheorem{definition}{Definicja}
\newtheorem{example}{Przykład}
\newtheorem{theorem}{Twierdzenie}
\newtheorem{lemma}{Lemat}
\newtheorem{remark}{Uwaga}
\newtheorem{proposition}{Stwierdzenie}

\title{Algorytmy zachłanne -- Kody Huffmana}
\author{}
\date{}

\begin{document}
\maketitle

\section{Kody Huffmana}

Często znaki są reprezentowane przy użyciu 8 bitów (np.\ ASCII).

Lepiej reprezentować znaki występujące często przy użyciu krótkich ciągów bitowych, a~znaki występujące rzadko przy użyciu długich sekwencji.

\begin{example}
W pliku jest 58 znaków, każdy znak jest elementem zbioru $\{a, e, i, s, t, r, n\}$.

\begin{center}
\begin{tabular}{|l|c|c|c|c|c|c|c|}
\hline
 & $a$ & $e$ & $i$ & $s$ & $t$ & $r$ & $n$ \\
\hline
liczba wystąpień (częstość) & 10 & 15 & 12 & 3 & 4 & 13 & 1 \\
\hline
Kod 1 (stałej długości) & 000 & 001 & 010 & 011 & 100 & 101 & 110 \\
\hline
Kod 2 (zmiennej dł.) & 001 & 01 & 10 & 00000 & 0001 & 11 & 00001 \\
\hline
Kod 3 (zmiennej dł.) & 001 & 01 & 10 & 0000 & 0001 & 11 & 000001 \\
\hline
\end{tabular}
\end{center}

Jeśli użyjemy kodu o~stałej długości słowa kodowego (Kod~1), to zakodowany plik ma $3 \times 58 = 174$ bitów.

Jeśli użyjemy zmiennej długości słowa kodowego (Kod~2), to zakodowany plik ma:
\[
    10 \cdot 3 + 15 \cdot 2 + 12 \cdot 2 + 3 \cdot 5 + 4 \cdot 4 + 13 \cdot 2 + 1 \cdot 5 = 146 \text{ bitów}
\]

Jeśli używamy kodu o~zmiennej długości słowa kodowego, to chcemy, aby spełniony był warunek, że żadne słowo kodowe nie jest prefiksem innego słowa kodowego.
\end{example}

\section{Kody prefiksowe}

Takie kody nazywamy \textbf{kodami prefiksowymi}.
% UWAGA: fragment w nawiasie o nazwie "bezprefiksowe" trudny do odczytu -- proszę sprawdzić.

Rozważmy Kod~3: ciąg bitów $0001001\ldots$

Nie możemy jednoznacznie odkodować:
\begin{itemize}
    \item $0001$ jest prefiksem $00010$,
    \item $0001 \to t$, $a$\ldots
    \item $00010 \to s$, $e$\ldots
\end{itemize}
Czyli słowo kodowe $0001$ jest prefiksem innego słowa kodowego -- nie jest to kod prefiksowy.

\subsection{Reprezentacja drzewowa}

Kody prefiksowe mają naturalną reprezentację drzewową. W~drzewie binarnym:
\begin{itemize}
    \item krawędź do lewego dziecka etykietujemy $0$,
    \item krawędź do prawego dziecka etykietujemy $1$.
\end{itemize}

\begin{center}
\begin{tikzpicture}[
    level distance=1.2cm,
    level 1/.style={sibling distance=3cm},
    level 2/.style={sibling distance=1.5cm},
    level 3/.style={sibling distance=1cm},
    every node/.style={circle, draw, minimum size=6mm, inner sep=1pt, font=\small}
]
\node {}
    child {node {}
        child {node {$a$} edge from parent node[left, draw=none] {0}}
        child {node {$e$} edge from parent node[right, draw=none] {1}}
        edge from parent node[left, draw=none] {0}
    }
    child {node {}
        child {node {$i$} edge from parent node[left, draw=none] {0}}
        child {node {}
            child {node {$s$} edge from parent node[left, draw=none] {0}}
            child {node {$t$} edge from parent node[right, draw=none] {1}}
            edge from parent node[right, draw=none] {1}
        }
        edge from parent node[right, draw=none] {1}
    };
\end{tikzpicture}

\smallskip
\textit{Przykładowe drzewo kodu (schematycznie).}
\end{center}

Słowo kodowe symbolu jest ścieżką bitów na ścieżce od korzenia do liścia.

\begin{remark}
Kod jest prefiksowy $\iff$ każdy symbol w~reprezentacji drzewowej jest w~liściu.
\end{remark}

\subsection{Etykiety w węzłach}

Etykietujemy wierzchołki w~następujący sposób: każdy wierzchołek $v$ dostaje etykietę, będącą sumą wystąpień znaków, które są reprezentowane przez liście poddrzewa o~korzeniu~$v$.

\section{Optymalny kod prefiksowy}

Niech $C$ -- alfabet -- zbiór znaków. Mamy $f \colon C \to \mathbb{N}$:

$f(c)$ -- częstość znaku $c$ z~alfabetu $C$.

$d_T(c)$ -- głębokość liścia reprezentującego znak $c$ w~drzewie $T$.

$d_T(c) = $ długość słowa kodowego znaku $c$ w~kodzie reprezentowanym przez drzewo $T$.

Liczba bitów w~zakodowanym pliku wynosi:
\[
    b_T = \sum_{c \in C} d_T(c) \cdot f(c)
\]

Kod prefiksowy reprezentowany przez drzewo $T$ nazywamy \textbf{optymalnym}, jeśli liczba $b_T$ jest najmniejsza możliwa.

\begin{definition}
Drzewo ukorzenione nazywamy \textbf{binarnym}, jeśli każdy węzeł ma co najwyżej dwójkę dzieci.
\end{definition}

\begin{definition}
Drzewo binarne nazywamy \textbf{regularnym}, jeśli każdy wierzchołek wewnętrzny ma dokładnie dwoje dzieci.
\end{definition}

\begin{proposition}
Drzewo binarne reprezentujące optymalny kod jest regularne.
\end{proposition}

\begin{proof}
Niech $T$ -- drzewo binarne reprezentujące optymalny kod. Załóżmy, że nie jest regularne. Wtedy istnieje wierzchołek wewnętrzny $v$, że ma dokładnie jedno dziecko $w$. Jeśli $v$ nie jest korzeniem, to niech $u$ -- rodzic wierzchołka $v$.

\begin{center}
\begin{tikzpicture}[
    every node/.style={circle, draw, minimum size=6mm, inner sep=1pt, font=\small},
    level distance=1cm, sibling distance=1.5cm
]
\node[draw=none] at (-1, 0.5) {$T$:};
\node {$u$}
    child {node {$v$}
        child {node {$w$}
            child {node {$a$}}
            child {node {$b$}}
        }
    }
    child {node {$\cdots$}};
\end{tikzpicture}
\hspace{3cm}
\begin{tikzpicture}[
    every node/.style={circle, draw, minimum size=6mm, inner sep=1pt, font=\small},
    level distance=1cm, sibling distance=1.5cm
]
\node[draw=none] at (-1, 0.5) {$T'$:};
\node {$u$}
    child {node {$w$}
        child {node {$a$}}
        child {node {$b$}}
    }
    child {node {$\cdots$}};
\end{tikzpicture}
\end{center}

Niech $T'$ będzie drzewem powstałym z~$T$ przez usunięcie $v$ i~połączenie $u$ z~$w$ (jeśli $v$ nie był korzeniem).

$b(T') < b(T)$, więc sprzeczność z~optymalnością.
\end{proof}

\section{Algorytm Huffmana}

Algorytm Huffmana -- konstruuje drzewo reprezentujące optymalny kod prefiksowy dla wejścia: alfabet $C$ i~funkcja częstości $f$.

Zaczynamy od $|C|$ wierzchołków izolowanych -- które stają się liśćmi, i~wykonujemy $|C| - 1$ operacji łączenia, zaczynając od najrzadszych znaków. One wtedy wylądują najgłębiej w~drzewie.

Będziemy używali kolejki priorytetowej z~kluczem $f(c)$.

Kolejkę priorytetową możemy zaimplementować używając \textbf{kopca binarnego}. Wtedy koszty operacji:

\begin{center}
\begin{tabular}{lc}
    \texttt{insert} & $O(\log n)$ \\
    \texttt{extract\_min} & $O(\log n)$ \\
    \texttt{create} & $O(n)$
\end{tabular}
\end{center}
gdzie $n = |C|$.

\bigskip

$z$ -- nowy wierzchołek; $\texttt{left}(z)$, $\texttt{right}(z)$ -- jego dzieci; $q$ -- kolejka priorytetowa.

\begin{algorithm}[H]
\caption{Algorytm Huffmana -- na wejściu $C$ i $f$}
\begin{algorithmic}[1]
\STATE $n \leftarrow |C|$ \hfill $O(n)$
\STATE Tworzymy $q$ na podstawie $f$ \hfill $O(n)$
\COMMENT{Mamy wykonać $n-1$ operacji merge}
\FOR{$i \leftarrow 1$ \textbf{to} $n - 1$}
    \STATE $x \leftarrow \texttt{left}(z) \leftarrow \texttt{extract\_min}(q)$ \hfill $O(\log n)$
    \STATE $y \leftarrow \texttt{right}(z) \leftarrow \texttt{extract\_min}(q)$ \hfill $O(\log n)$
    \STATE $f(z) \leftarrow f(x) + f(y)$
    \STATE $\texttt{insert}(z)$ \hfill $O(\log n)$
\ENDFOR
\end{algorithmic}
\end{algorithm}

Algorytm konstruuje drzewo binarne, ma $n - 1$ operacji mergeowania i~każdy wierzchołek wewnętrzny ma 2 dzieci.

\subsection{Przykład}

Krok 0 -- kolejka: $(h, 1)$, $(s, 3)$, $(t, 4)$, $(a, 10)$, $(i, 12)$, $(r, 13)$, $(e, 15)$.

\textbf{Krok 1:} Łączymy $h{:}1$ i $s{:}3 \to$ nowy węzeł $4$.

\begin{center}
\begin{tikzpicture}[
    every node/.style={circle, draw, minimum size=6mm, inner sep=1pt, font=\small},
    level distance=1cm, sibling distance=1.5cm
]
\node {4}
    child {node {$h$}}
    child {node {$s$}};
\end{tikzpicture}
\end{center}

Kolejka: $(4)$, $(t, 4)$, $(a, 10)$, $(i, 12)$, $(r, 13)$, $(e, 15)$.

\textbf{Krok 2:} Łączymy $4$ i $t{:}4 \to$ nowy węzeł $8$.

\begin{center}
\begin{tikzpicture}[
    every node/.style={circle, draw, minimum size=6mm, inner sep=1pt, font=\small},
    level distance=1cm, sibling distance=1.8cm
]
\node {8}
    child {node {4}
        child {node {$h$}}
        child {node {$s$}}
    }
    child {node {$t$}};
\end{tikzpicture}
\end{center}

Kolejka: $(8)$, $(a, 10)$, $(i, 12)$, $(r, 13)$, $(e, 15)$.

\textbf{Krok 3:} Łączymy $8$ i $a{:}10 \to$ nowy węzeł $18$.

\begin{center}
\begin{tikzpicture}[
    every node/.style={circle, draw, minimum size=6mm, inner sep=1pt, font=\small},
    level distance=1cm, sibling distance=2.2cm,
    level 2/.style={sibling distance=1.2cm}
]
\node {18}
    child {node {8}
        child {node {4}
            child {node {$h$}}
            child {node {$s$}}
        }
        child {node {$t$}}
    }
    child {node {$a$}};
\end{tikzpicture}
\end{center}

Kolejka: $(i, 12)$, $(r, 13)$, $(e, 15)$, $(18)$.

\textbf{Krok 4:} Łączymy $i{:}12$ i $r{:}13 \to$ nowy węzeł $25$.

\begin{center}
\begin{tikzpicture}[
    every node/.style={circle, draw, minimum size=6mm, inner sep=1pt, font=\small},
    level distance=1cm, sibling distance=1.5cm
]
\node {25}
    child {node {$i$}}
    child {node {$r$}};
\end{tikzpicture}
\end{center}

Kolejka: $(e, 15)$, $(18)$, $(25)$.

\textbf{Krok 5:} Łączymy $e{:}15$ i $18 \to$ nowy węzeł $33$.

\begin{center}
\begin{tikzpicture}[
    every node/.style={circle, draw, minimum size=6mm, inner sep=1pt, font=\small},
    level distance=1cm, sibling distance=2.5cm,
    level 2/.style={sibling distance=1.5cm},
    level 3/.style={sibling distance=1cm}
]
\node {33}
    child {node {$e$}}
    child {node {18}
        child {node {8}
            child {node {4}
                child {node {$h$}}
                child {node {$s$}}
            }
            child {node {$t$}}
        }
        child {node {$a$}}
    };
\end{tikzpicture}
\end{center}

Kolejka: $(25)$, $(33)$.

\textbf{Krok 6:} Łączymy $25$ i $33 \to$ nowy węzeł $58$ (korzeń).

\begin{center}
\begin{tikzpicture}[
    every node/.style={circle, draw, minimum size=6mm, inner sep=1pt, font=\small},
    level distance=1cm, sibling distance=4cm,
    level 2/.style={sibling distance=2.5cm},
    level 3/.style={sibling distance=1.8cm},
    level 4/.style={sibling distance=1.2cm}
]
\node {58}
    child {node {25}
        child {node {$i$}}
        child {node {$r$}}
    }
    child {node {33}
        child {node {$e$}}
        child {node {18}
            child {node {8}
                child {node {4}
                    child {node {$h$}}
                    child {node {$s$}}
                }
                child {node {$t$}}
            }
            child {node {$a$}}
        }
    };
\end{tikzpicture}
\end{center}

Odczytując kody z~drzewa (lewo $= 0$, prawo $= 1$):

\begin{center}
\begin{tabular}{cl}
    $i$: & 00 \\
    $r$: & 01 \\
    $e$: & 10 \\
    $a$: & 111 \\
    $t$: & 1101 \\
    $h$: & 11000 \\
    $s$: & 11001 \\
\end{tabular}
\end{center}

\textcolor{red}{\textbf{UWAGA:} Kody dla $h$ i $s$ mogą być zamienione -- trudne do odczytu z~notatek. Proszę zweryfikować.}

\subsection{Analiza złożoności}

Złożoność czasowa wynosi $O(n \log n)$:
\begin{itemize}
    \item Inicjalizacja (linie 1--2): $O(n)$
    \item Pętla wykonuje się $n - 1$ razy, każda iteracja: $O(\log n)$
    \item Razem: $O(n) + (n-1) \cdot O(\log n) = O(n \log n)$
\end{itemize}

\end{document}
